\documentclass[10pt,a4paper]{article}
\usepackage[margin=18mm,headheight=15pt]{geometry}
\usepackage{amsmath,booktabs,array,makecell,xcolor,graphicx,fancyhdr,draftwatermark,hyperref}
\usepackage[T1]{fontenc}
\usepackage{lmodern}
\hypersetup{pdftitle={Experiment Matrix GUIDE: Prospective Targets, Not Observed Results},pdfsubject={Internal planning scenario; not preregistered; not for submission},pdfauthor={Earning Roles project},colorlinks=true,urlcolor=blue!60!black}
\SetWatermarkText{GUIDE --- NOT OBSERVED}
\SetWatermarkColor{red!14}
\SetWatermarkAngle{35}
\SetWatermarkScale{0.70}
\pagestyle{fancy}\fancyhf{}
\fancyhead[L]{\color{red!70!black}\bfseries GUIDE: PROSPECTIVE TARGETS}
\fancyhead[R]{\color{red!70!black}NOT EXPERIMENTAL RESULTS}
\fancyfoot[L]{\color{red!70!black}v1 / 2026-10-08 / NOT PREREGISTERED / NOT FOR SUBMISSION}
\fancyfoot[R]{\thepage}
\setlength{\parindent}{0pt}\setlength{\parskip}{5pt}
\setlength{\tabcolsep}{4pt}\renewcommand{\arraystretch}{1.13}
\newcommand{\rare}{\textsc{rare}}
\newcommand{\metric}[1]{\makecell[c]{#1}}
\newcommand{\status}{\textcolor{red!70!black}{\textbf{Guide numbers are assumptions or goals, never measurements.}}}
\newcommand{\lead}[1]{\vspace{4pt}\textbf{#1}\par}
\begin{document}
{\Large\bfseries Experiment Matrix Guide}\hfill{\small Earning Roles}
\par\status
This independent guide mirrors the three proposed body result tables. S1 is a
\emph{working target scenario}, not a fitted forecast. Comparator rows are
planning reference points; RARE rows are desired endpoints. They neither freeze
the experiment cards nor replace qualified roots, baselines, labels or costs.
All $n$ entries are \emph{planned streams}, with zero guide experiments executed.

\lead{Measured anchors: one PIPE3 development run, original v2 scorers}
\begin{center}\small
\begin{tabular}{@{}lcccc@{}}\toprule
Policy & Producer checks & Recipient checks & Adoption checks & One update (ms)\\\midrule
No update & 2/3 & 3/3 & 1/2 & Not run\\
Ctx.-trust linear & 2/3 & 3/3 & 1/2 & 0.520\\
RARE & 2/3 & 3/3 & 1/2 & 0.297\\\bottomrule
\end{tabular}\end{center}
Eight real API calls completed. All three target choices were the same peer.
Updates had no subsequently executed target. Check fractions are not task
accuracy, and these rows are not independent trials. Thus no measured quality
advantage is used to construct S1. Serialized state was 3,408 bytes for the
linear control and 4,688 bytes for RARE; this is not resident memory.

\lead{Table 1 counterpart: ArtifactRole prediction and future allocation --- S1 targets}
\begin{center}\small
\begin{tabular}{@{}lrrrrrr@{}}\toprule
Method & Brier $\downarrow$ & Calib. $\downarrow$ & $U\uparrow$ & $\Delta A$ & Cost & $n^{\rm plan}$\\\midrule
Uniform & 0.250 & 0.000 & 0.400 & 0.00 & 1.00 & 18\\
No update & 0.250 & 0.000 & 0.400 & 0.00 & 1.00 & 18\\
Raw acceptance & 0.233 & 0.150 & 0.398 & 0.06 & 1.02 & 18\\
Terminal only & 0.220 & 0.100 & 0.419 & 0.10 & 1.01 & 18\\
Ctx.-trust linear & 0.215 & 0.070 & 0.447 & 0.15 & 1.03 & 18\\
Pooled controller & 0.215 & 0.070 & 0.447 & 0.18 & 1.03 & 18\\
Meta-Team L2-public$^{\dagger}$ & 0.215 & 0.070 & 0.447 & 0.15 & 1.03 & 18\\
RARE & 0.212 & 0.040 & 0.495 & 0.20 & 1.05 & 18\\
\bottomrule\end{tabular}\end{center}
\textbf{GUIDE TARGET SCENARIO; NOT OBSERVED.} $^{\dagger}$The Meta-Team adapter is
unqualified: its row is a conditional strong-control proxy, \emph{not} an
estimate of the published method. RARE is allowed to cost 5\% more than no
update and is not assumed to win every column. A greater $\Delta A$ is not
automatically beneficial.

\lead{Definitions and reproducible calculation}
Let $Q$ be future independent terminal success, including full recipient use;
$C$ is mean \emph{complete} per-assignment cost divided by no-update cost on the
same root and stream. In this guide only, $U=Q-0.10C$. S1 assumes
$Q_0=.50$, $Q_{\rm ctx}=.55$, $Q_{\rm RARE}=.60$. The measured 1/2 adoption
check fraction merely warns against assuming near-perfect quality; it is not
a fitted estimate of $Q$. Then $U_{\rm RARE}=.60-.10(1.05)=.495$ and
$U_{\rm ctx}=.55-.10(1.03)=.447$: the \emph{desired} difference is .048.

For prediction, two equally weighted, evaluation-only strata have assumed true
rates $q=(.30,.70)$. With forecasts $p$, expected Brier is
$\frac12\sum_{g=1}^2[q_g(1-q_g)+(p_g-q_g)^2]$.
RARE's target $p=(.26,.74)$ gives $.21+.04^2=.2116$.
Calibration is population ECE in fixed forecast bins $[0,.5)$ and $[.5,1]$.
Constant .5 forecasts have ECE zero but Brier .25, so calibration alone does
not show informative judgment. True rates are never selector inputs.

\clearpage
\lead{Table 2 counterpart: PeerSelect service and drift --- targets, not forecasts}
\status
\begin{center}\small
\begin{tabular}{@{}lrrrrrrr@{}}\toprule
Method & \metric{Update p95\\(ms)} & Backlog & \metric{State\\(KiB)} & Payoff & Forget & \metric{Recovery\\(steps)} & $n^{\rm plan}$\\\midrule
Local uniform & -- & 0 & $\leq$8 & 0.50 & $\leq$0.00 & $>50$ & 18\\
No history & -- & 0 & $\leq$8 & 0.50 & $\leq$0.00 & $>50$ & 18\\
History only & $\leq$5 & $\leq$4 & $\leq$64 & 0.54 & $\leq$0.02 & $\leq$40 & 18\\
Graph-IPD reference$^{\dagger}$ & $\leq$5 & $\leq$4 & $\leq$1024 & 0.56 & $\leq$0.02 & $\leq$20 & 18\\
RARE & $\leq$5 & $\leq$4 & $\leq$64 & 0.56 & $\leq$0.01 & $\leq$20 & 18\\
\bottomrule\end{tabular}\end{center}
\textbf{GUIDE DESIGN ENVELOPES; NOT OBSERVED.} $^{\dagger}$Graph-IPD values are
conditional comparison targets, not its measured speed or adaptation. Its
upstream smoke reward 14.15 used a different aggregation and cannot be converted
into the normalized payoff here. Memory limits are engineering allowances,
not estimates of the reference implementation. RARE and the strong reference
are deliberately assigned the same payoff and recovery targets.

\lead{Service workload and arithmetic}
These budgets cover a CPU selector with cached 64-dimensional features, queue
handling, legal small-head update, state serialization and logging. The .297 ms
and .520 ms measured anchors cover only one existing update; 5 ms is a new
service allowance, not a p95 extrapolation. Neural encoding, retrieval/model
training and end-to-end publish/read service must be measured separately before
any full real-time-training claim. Backbone and updater remain open.

The guide workload is four feedback arrivals every 40 ms, giving
$\lambda=4/.040=100$ events/s. A \emph{maximum} 5 ms service envelope, stronger
than p95 alone, processes a burst in $4\times5=20$ ms before the next burst.
The backlog $\leq4$ target in the table is conditional on this \emph{maximum}
5 ms service bound; it includes the event in service. Actual p95 cannot by
itself guarantee backlog; both are test targets,
and burst tails, disk stalls and asynchronous corrections must be recorded.

The RARE dynamic-state allowance is
\[
 B=4096+4(64)(8)+128(256)=38912\ \text{bytes}=38.0\ \text{KiB}<64\ \text{KiB}.
\]
This reserves four float64 head vectors, 128 retained 256-byte records and
4 KiB of fixed metadata. It is a proposed representation budget, not a statement
that the current model uses it. Count serialized metadata and every retained
feedback/replay record; model weights, full ledger, runtime resident memory and
external artifact storage are reported separately, not silently discarded.
No semantic compression quality is inferred from a byte limit.

\lead{Payoff, forgetting and recovery}
For the guide only, use the IPD payoff matrix $T=5,R=3,P=1,S=0$ and divide
selected-pair payoff by 5. Values .50 and .56 mean raw payoff 2.50 and 2.80,
within $[0,5]$. Only selected realized payoff may enter learning; unchosen
oracle outcomes, if available for evaluation, never enter feedback.

Forgetting is $\max(0,Q_{\rm old,pre}-Q_{\rm old,post})$ on a fixed old-task
holdout. A target .01 is one percentage point, not 1\% relative loss. Recovery
is the earliest post-drift step whose independent evaluation reaches at least
90\% of the fixed post-drift reference payoff for three consecutive windows.
The reference is frozen before the comparison and costs are matched. ``$>50$''
for fixed policies is the guide scenario of drift that requires learning,
not a prediction that every drift defeats a fixed policy. The exact evaluation
window, reference and confidence rule remain to be frozen in the real card.

\lead{Why these targets are restrained}
Sub-millisecond core updates make a 5 ms cached-feature allowance worth testing;
they say nothing about an unmeasured backbone forward/backward pass. No measured
drift or forgetting data exist. Therefore the adaptation cells are desired
constraints rather than extrapolated expectations, and strong controls can
equal or outperform RARE when tested.

\clearpage
\lead{Table 3 counterpart: mechanisms and responsibility --- S1 targets}
\status
\begin{center}\small
\begin{tabular}{@{}lrrrr@{}}\toprule
\multicolumn{5}{l}{\textbf{(a) Mechanism ablations}}\\
Variant & $U$ & $\Delta A$ & False attrib. & Cost\\\midrule
Full RARE & 0.495 & 0.20 & 0.000 & 1.05\\
No public evidence & 0.428 & 0.10 & 0.000 & 1.02\\
No delayed credit & 0.447 & 0.15 & 0.000 & 1.03\\
Neither evidence nor credit & 0.400 & 0.00 & 0.000 & 1.00\\
No ownership gate & 0.435 & 0.24 & 0.138 & 1.05\\
No locality & 0.492 & 0.22 & 0.000 & 1.18\\
No selected-only rule & \multicolumn{4}{l}{\textit{N/A: feedback contract not defined}}\\
No correction / replay & 0.466 & 0.20 & 0.000 & 1.04\\
\midrule
\multicolumn{5}{l}{\textbf{(b) False producer-credit rate in negative-ownership stress}}\\
Variant & Recipient-owned & Mixed & Unobserved & Resource failure\\\midrule
Full RARE & 0.000 & 0.000 & 0.000 & 0.000\\
No ownership gate & 0.250 & 0.150 & 0.100 & 0.050\\
\bottomrule\end{tabular}\end{center}
\textbf{GUIDE TARGET SCENARIO; NOT OBSERVED.} The selected-only ablation cannot
be assigned a legal number under the current observed-feedback contract. It
requires a separately specified diagnostic intervention; pretending unchosen
labels are available would make the table mathematically invalid.

\lead{Four-cell consistency and testable interaction}
In S1, let evidence publication and delayed credit be independently enabled:
\[
\begin{array}{c|cc}
 & \text{No credit} & \text{Credit}\\\hline
\text{No evidence} & U_{00}=.400 & U_{01}=.428\\
\text{Evidence} & U_{10}=.447 & U_{11}=.495
\end{array}
\]
Then the publication-only increment is .047, the credit-only increment .028,
and the interaction is
\[
 I=U_{11}-U_{10}-U_{01}+U_{00}=.020;
 \qquad .047+.028+.020=.095=U_{11}-U_{00}.
\]
This is the interaction we \emph{aim to test}, not evidence that it exists.
Quality inputs are respectively .50, .53, .55 and .60; cost inputs are 1.00,
1.02, 1.03 and 1.05. Their utility cells use the same coefficient as Table 1.

The evidence switch controls the \emph{selector's public read}, not existence of
the operator's immutable assignment/provenance receipt. In the credit-only cell,
that receipt remains available for legal post-target selected-only credit, while
its source observation is withheld from the selector. All cells retain the same
annotation and compute budget and reward contract; each arm uses its own realized
future outcome and history. The credit switch controls only parameter updates.
This separation requires runner qualification before the factorial comparison;
without it the interaction is not identifiable.

An assignment-change rate is the fraction of paired decisions whose chosen peer
differs from no update using the same menu and random draw. Under a quality-only
selection-effect model, $|Q-Q_0|\leq\Delta A$; every S1 row satisfies this bound.
For RARE, a .10 quality gain with .20 changed choices requires an average .50
gain on changed decisions. This is a demanding target, not a free consequence
of changing probabilities. Real policies may also change delivery history;
that pathway must be measured rather than forced into this simple model.

\lead{Safety denominators and an explicit trade-off}
False attribution means producer reward/penalty issued in a case whose frozen
contract and independent evidence prohibit producer credit, divided by all
tested negative-ownership cases. UNKNOWN and failures stay in the denominator
with distinct status. Full RARE's zero is a fail-closed invariant target, not
an observed zero-error estimate. No-gate values (.25,.15,.10,.05) are stress
planning assumptions. With equally many cases per category, their overall rate
is $(.25+.15+.10+.05)/4=.1375$, matching panel (a).

No locality may slightly improve quality (.61 versus .60) but increase complete
cost to 1.18, giving $U=.492$ versus .495. This retains a quality/cost trade-off
instead of forcing every ablation to lose on every metric. A future result can
overturn any S1 ordering; changing the guide never changes recorded results.

\clearpage
\lead{Sensitivity, precision and execution scale}
\status
\begin{center}\small
\begin{tabular}{@{}lrrrrr@{}}\toprule
Scenario & $Q_R$ & $C_R$ & $U_R$ & $U_{\rm ctx}$ & Difference\\\midrule
Null: no quality gain & 0.50 & 1.05 & 0.395 & 0.397 & -0.002\\
Small gain & 0.56 & 1.05 & 0.455 & 0.447 & 0.008\\
S1: working target & 0.60 & 1.05 & 0.495 & 0.447 & 0.048\\
Cost overrun & 0.60 & 1.20 & 0.480 & 0.447 & 0.033\\
Strong control wins & 0.53 & 1.05 & 0.425 & 0.497 & -0.072\\
\bottomrule\end{tabular}\end{center}
These scenarios intentionally include a negative result, a near-tie and a strong
control win. They
are not a probability distribution. Unknown quality effects have no fitted
posterior or prediction interval. No invented 95\% result intervals or
significance marks appear in the target matrices.

\lead{What the target difference would require}
For the same S1 $\Delta Q=.05$ and $\Delta C=.02$,
$\Delta U=.05-\lambda(.02)$. At guide sensitivity weights
$\lambda=.05,.10,.20$, the differences are .049, .048 and .046.
With no quality gain they are negative (.001, .002 and .004 in magnitude).
The actual complete-cost conversion, utility coefficient and threshold must be
frozen before confirmation. Wall time is not a substitute for money, GPU,
tokens, tools, human effort and retry costs.

Assuming, not observing, a paired-stream standard deviation $s_D=.06$ and
target difference .048, the two-sided normal approximation for 80\% power is
\[
 n\approx\left(\frac{(1.96+.8416).06}{.048}\right)^2,
 \quad \lceil n\rceil=13.
\]
The illustrative plan uses 18 paired stream differences for one primary contrast
(six per each of at least three qualified structural roots). With
$t_{.975,17}=2.1098$, the assumed
95\% half-width is $2.1098(.06)/\sqrt{18}=0.030$.
This is a \emph{precision calculation}, not a measured interval or final sample
size. It describes conditional within-root precision under the assumed SD; it
does not establish 80\% power for transfer across unseen structural roots.
Root hierarchy, multiple primary comparisons and larger pilot variance can
require more streams. Freeze the root-aware analysis and primary contrast, then
re-estimate sample size from a development pilot. A stream contains many episodes; they are never
counted as independent streams. Root split and seeds are not selected here.

For an independent zero-error safety diagnostic, 300 negative cases would give
the one-sided 95\% upper bound $1-.05^{1/300}\approx.00994$.
Repeated cases within streams may be dependent; then this Bernoulli calculation
does not apply and the registered clustered analysis must be used.

\lead{Budget reality check before any launch}
Eighteen streams, eight ArtifactRole arms and 30 future assignments per stream
imply 4,320 target episodes. At three new calls per fresh target and three
initial-source calls per arm/stream, the skeleton is 13,392 calls.
The eight-call measured run averaged 1,814.375 input-plus-output tokens per call;
holding that average gives approximately 24.30 million tokens and
40.9 summed API hours. A 2$\times$ planning multiplier gives roughly twice
these totals. These are workload projections, not full budgets: profile calls,
retries, tools, scorer CPU, communication, model training and overhead must be
added explicitly. No pricing or parallel wall-clock guarantee is inferred.

This is a possible \emph{confirmation-scale} target, not the next experiment.
First qualify roots, label/update legality and same-information baselines, then
use a bounded development pilot to estimate stream variance and whether the
signal exists. Freeze the complete cards, precision and stopping rules before
confirmation. A null or reversed strong-control comparison requires analyzing
the recorded failure and revising the method; it cannot be repaired by replacing
the measured numbers or lowering the Goal. This guide does not alter execution
budgets or launch resources; no experiment was performed to fill its targets.

\lead{Version and source discipline}
The PDF is independent of the canonical paper. Its editable source, target
specification, arithmetic checks and raw-source hash are archived with v1.
Every page carries a visible watermark and status header. The measured data
remain unchanged; guide cells never feed the paper's real result tables.
The active Goal and three acceptance standards remain authoritative.
\end{document}
